\section{ファンアウト1の区別対による下界}
\label{app:fanout-one-lower}
関連CFG論文の括弧付き倍増言語族を用いる．
同言語族がClark--Eyraud代入可能であることのみ
TCSへの提出改訂稿~\cite[Proposition~6.4(i) and Appendix~C]{kuriyamaCFG}
から引用し，下界は独立に証明する．
\begin{lemma}[区別対]
\label{lem:distinguishing-pair}
集合駆動型作用素$\mathcal O$について
$L'\subsetneq L$の両方に特徴標本が存在するとき，
$L$の各特徴標本は$L\setminus L'$と交わる．
\end{lemma}
\begin{proof}
$L$の特徴標本$F$が$F\subseteq L'$と仮定する．
$L'$の特徴標本$F'$を取って$K=F\cup F'$と置くと，
$F\subseteq K\subseteq L$から
$L(\mathcal O(K))=L$，
$F'\subseteq K\subseteq L'$から
$L(\mathcal O(K))=L'$が従い，矛盾する．
\end{proof}

$\Gamma_{\rm br}=\{\mathtt a,\lfloor,\rfloor,\lceil,\rceil\}$，
$U_1=\{1,\zeta\}$，
$1\zeta=\zeta1=\zeta^2=\zeta$とする．
$s_{\mathrm U}:\Gamma_{\rm br}^+\twoheadrightarrow U_1$は
$\mathtt a,\lfloor,\rfloor$を$1$，
$\lceil,\rceil$を$\zeta$に送る．
$G_n$を次のファンアウト1のMCFGとする．
\[
 S\to E_n,\quad E_0\to\mathtt a,\quad
 E_i\to\lfloor E_{i-1}E_{i-1}\rfloor
 \mid\lceil E_{i-1}\rceil\quad(1\le i\le n).
\]
$L_n=L(G_n)$とし，
$w_0^*=\mathtt a$，
$w_i^*=\lfloor w_{i-1}^*w_{i-1}^*\rfloor$，
$L_n^-=L_n\setminus\{w_n^*\}$と置く．
\begin{proposition}[倍増言語族の性質]
\label{prop:fanout-one-family}
各$n\ge1$について，
\begin{enumerate}[label=\textup{(\roman*)},leftmargin=*]
\item $G_n$は縮約済み良形，規則ランク2，
$\|G_n\|=O(n)$，$\theta_{G_n}\le4n$である．
\item $w_n^*$は$L_n$の中で天井括弧を含まない唯一の語で，
$|w_n^*|=3\cdot2^n-2$である．
特に$L_n^-=L_n\cap\widehat{s}_{\mathrm U}^{-1}(\{\zeta\})$．
\item $L_n,L_n^-\in\Cmcf{1}{s_{\mathrm U}}$である．
\end{enumerate}
\end{proposition}
\begin{proof}
(i) 全非終端記号は到達可能かつ生産的である．
各子出現は一度だけ使用され，ファンアウト1なので
成分の置換も併合も起こらない．
規則数は$2n+2$であり各規則は定数長である．
選んだ規則以外をすべて$\lceil\cdot\rceil$型とする．
分岐規則$E_i\to\lfloor E_{i-1}E_{i-1}\rfloor$の証人は
$u=\lceil^{i-1}\mathtt a\rceil^{i-1}$として
$\lceil^{n-i}\lfloor uu\rfloor\rceil^{n-i}$を取ればよく，
長さは$2(n-i)+2+2(2i-1)=2n+2i\le4n$．
他の規則は長さ$2n+1$の証人で実現できる．
(ii) 天井括弧を避ける導出は完全二分木導出に限られる．
$|w_0^*|=1$，$|w_i^*|=2|w_{i-1}^*|+2$より
指定した長さ公式を得る．
(iii) $L_n$は関連CFG改訂稿によって
Clark--Eyraud代入可能であり，型を細分した
$s_{\mathrm U}$の下でも代入可能である．
型繊維$\widehat s_{\mathrm U}^{-1}(\{\zeta\})$との交差は，
同じ型の断片の置換について代入可能性を保つので，
$L_n^-$にも成り立つ．両者は有限だからCFLである．
\end{proof}
